\documentclass[12pt]{article}
\usepackage[margin=1in]{geometry}
\usepackage{amsmath, amssymb, graphicx, booktabs, subcaption,url}

\usepackage{setspace}
\usepackage{url}

\doublespacing

\usepackage[round]{natbib}
\usepackage{url}
\usepackage{hyperref}
\providecommand{\natexlab}[1]{#1}
\providecommand{\doi}[1]{doi: \url{#1}}  % if doi is used

% Manual definitions (in case BibTeX inserts these macros)


\bibliographystyle{plainnat}
\graphicspath{{../Figs/}}
\title{Rental Household Size as a Novel Measure of Demand: Evidence from the Rental Density Index (RDI)}

\author{
	Matt Larriva$^{1}$ \\ \texttt{matt.larriva@brookfield.com}
}

\date{\vspace{-5ex}}  % Remove date spacing if needed

\begin{document}
	
	\maketitle
	
	\begin{center}
		$^{1}$Department of Real Estate Investments, Brookfield Asset Management, New York, USA
	\end{center}
	
	\vspace{1em}
	
	\section*{Abstract}
	Traditional measures of rental housing demand, such as occupancy and absorption, are constrained by supply and fail to capture latent demand pressures in fully leased markets. This paper introduces the \emph{Rental Density Index (RDI)}, a novel metric defined as the number of occupants per rental unit, to measure underlying crowding in rental households. Using data from the 200 largest U.S. metropolitan areas from 2005 to 2023, we show that changes in RDI predict future rent growth, outperforming standard time-series models. We employ a two-stage least squares approach that instruments for RDI growth using local birth cohorts from 20 years prior to identify a causal relationship. Markets with rising RDI see lower near-term rent growth but are associated with future upward pressure. Forecasting tests confirm that RDI-based models outperform ARIMA and naïve approaches across multiple horizons. These findings establish RDI as a scalable, behaviorally grounded indicator of rental housing demand and a practical tool for investors and policymakers.\cite{betancourt2022us}
	
	\vspace{1em}
	\noindent\textbf{Keywords:} density, latent demand, multifamily housing, forecasting, rent growth
	
	\newpage
\section{Introduction}
\newpage
\bibliography{mwe}


\end{document}
